\documentclass[11pt]{article}

% Change "review" to "final" to generate the final (sometimes called camera-ready) version.
% Change to "preprint" to generate a non-anonymous version with page numbers.
\usepackage[review]{acl}
\usepackage{float}
% Standard package includes
\usepackage{times}
\usepackage{latexsym}
\usepackage{booktabs}
% For proper rendering and hyphenation of words containing Latin characters (including in bib files)
\usepackage[T1]{fontenc}
% For Vietnamese characters
% \usepackage[T5]{fontenc}
% See https://www.latex-project.org/help/documentation/encguide.pdf for other character sets
\usepackage{svg}
\usepackage{amsmath}
% This assumes your files are encoded as UTF8
\usepackage[utf8]{inputenc}

% This is not strictly necessary, and may be commented out,
% but it will improve the layout of the manuscript,
% and will typically save some space.
\usepackage{microtype}
\usepackage{multirow}

% This is also not strictly necessary, and may be commented out.
% However, it will improve the aesthetics of text in
% the typewriter font.
\usepackage{inconsolata}

%Including images in your LaTeX document requires adding
%additional package(s)
\usepackage{graphicx}
\usepackage{svg}
% If the title and author information does not fit in the area allocated, uncomment the following
%
%\setlength\titlebox{<dim>}
%
% and set <dim> to something 5cm or larger.

\title{Q-Bimodal Tree Attention: Fixed-Cardinality Dispatch for Efficient Tree-Structured LLM Decoding}

% Author information can be set in various styles:
% For several authors from the same institution:
% \author{Author 1 \and ... \and Author n \\
%         Address line \\ ... \\ Address line}
% if the names do not fit well on one line use
%         Author 1 \\ {\bf Author 2} \\ ... \\ {\bf Author n} \\
% For authors from different institutions:
% \author{Author 1 \\ Address line \\  ... \\ Address line
%         \And  ... \And
%         Author n \\ Address line \\ ... \\ Address line}
% To start a separate ``row'' of authors use \AND, as in
% \author{Author 1 \\ Address line \\  ... \\ Address line
%         \AND
%         Author 2 \\ Address line \\ ... \\ Address line \And
%         Author 3 \\ Address line \\ ... \\ Address line}

\author{First Author \\
  Affiliation / Address line 1 \\
  Affiliation / Address line 2 \\
  Affiliation / Address line 3 \\
  \texttt{email@domain} \\\And
  Second Author \\
  Affiliation / Address line 1 \\
  Affiliation / Address line 2 \\
  Affiliation / Address line 3 \\
  \texttt{email@domain} \\}

%\author{
%  \textbf{First Author\textsuperscript{1}},
%  \textbf{Second Author\textsuperscript{1,2}},
%  \textbf{Third T. Author\textsuperscript{1}},
%  \textbf{Fourth Author\textsuperscript{1}},
%\\
%  \textbf{Fifth Author\textsuperscript{1,2}},
%  \textbf{Sixth Author\textsuperscript{1}},
%  \textbf{Seventh Author\textsuperscript{1}},
%  \textbf{Eighth Author \textsuperscript{1,2,3,4}},
%\\
%  \textbf{Ninth Author\textsuperscript{1}},
%  \textbf{Tenth Author\textsuperscript{1}},
%  \textbf{Eleventh E. Author\textsuperscript{1,2,3,4,5}},
%  \textbf{Twelfth Author\textsuperscript{1}},
%\\
%  \textbf{Thirteenth Author\textsuperscript{3}},
%  \textbf{Fourteenth F. Author\textsuperscript{2,4}},
%  \textbf{Fifteenth Author\textsuperscript{1}},
%  \textbf{Sixteenth Author\textsuperscript{1}},
%\\
%  \textbf{Seventeenth S. Author\textsuperscript{4,5}},
%  \textbf{Eighteenth Author\textsuperscript{3,4}},
%  \textbf{Nineteenth N. Author\textsuperscript{2,5}},
%  \textbf{Twentieth Author\textsuperscript{1}}
%\\
%\\
%  \textsuperscript{1}Affiliation 1,
%  \textsuperscript{2}Affiliation 2,
%  \textsuperscript{3}Affiliation 3,
%  \textsuperscript{4}Affiliation 4,
%  \textsuperscript{5}Affiliation 5
%\\
%  \small{
%    \textbf{Correspondence:} \href{mailto:email@domain}{email@domain}
%  }
%}

\begin{document}
\maketitle
\begin{abstract}
Tree-structured decoding workloads generate many continuations from a shared context, causing redundant KV-prefix reads and inefficient attention execution. We observe that tree attention is query-bimodal: shared prefixes favor wide prefill-style kernels, while branch-specific tails favor decode-native execution. We propose a fixed-cardinality dispatcher that dynamically selects among a small set of shared-prefix and decode-tail strategies using a calibrated cost model. Keeping the dispatch space small while adapting to prefix length, batch size, and tail shape. On H100 GPUs with Llama-3.2-1B, Llama-3.1-8B and Llama-3-70B, our method consistently outperforms paged decoding, fixed shared-prefix execution, and existing dynamic tree-attention planning, achieving up to 6.4× over PagedAttention and up to 2.73× over the strongest competing baseline.
\end{abstract}

\section{Introduction}

Tree-structured generation is a common pattern in modern LLM inference.
Beam search \cite{vijayakumar2018diversebeamsearchdecoding}, speculative decoding \cite{10.1145/3620666.3651335, li2025eagle3}, few-shot
learning \cite{10.5555/3495724.3495883}, multi-step reasoning \cite{yao2023treethoughtsdeliberateproblem, hao2023reasoninglanguagemodelplanning} all generate multiple continuations from a
shared context. Although these workloads differ in how branches are produced,
they induce the same attention-time structure: a long shared prefix that fans out
into many branch-specific continuations. If each branch is served as an
independent request, every branch repeatedly rereads the same shared KV prefix
from HBM, even though the prefix is identical across branches.

Prior systems reduce this redundancy in different ways, as summarized in
Table~\ref{method-comparison}, where $K$ denotes the number of leaves in the
tree. PagedAttention \cite{10.1145/3600006.3613165} improves KV-cache
management, but does not share prefix reads across branches. DeFT
\cite{yao2025deft} and Cascade Attention \cite{ye2025flashinfer} reuse prefixes
across multiple KV spans, but their execution structure is fixed rather than
selected dynamically for each decoding step. FastTree \cite{pan2025fasttree}
dynamically plans over tree-structured KV caches, but its search space grows
with the input tree. Moreover, it executes all levels using a prefill-style
kernel template, including branch-specific tails that are better suited to
decode-style execution. As shown in Figure~\ref{fig:poor_query_padding}, this
design forces tail queries to use the same tile size as shared-prefix queries.
Because GPU kernels require fixed tile sizes, the small query counts in tail
regions lead to substantial query-tile padding waste.

We observe that tree attention is query-bimodal: shared-prefix attention favors
wide prefill-style execution, while branch-specific tail attention favors
decode-native execution. To address this mismatch, we propose Q-Bimodal Tree
Attention, which dynamically analyzes the tree structure to select among a small
set of execution strategies and launches a dedicated decode kernel when doing so
is more efficient.


\begin{figure}
    \centering
    \includegraphics[width=0.9\linewidth]{padding_tree.png}
    \caption{Poor query padding for single kernel tree attentions. }
    \label{fig:poor_query_padding}
\end{figure}

Our contributions are as follows:
\begin{itemize}
    \item We identify query-shape bimodality as a central bottleneck in
    tree-structured attention: shared prefixes require wide query tiles, while
    branch-specific tails require decode-native execution.
    \item We show that a dedicated decoding kernel can remove tail-side
    query-tile padding, but that it is not universally faster because
    it introduces launch, merge, and occupancy overheads.
    \item We design a tree-attention kernel with query-bimodal runtime
    dispatch, achieving best performance across diverse decoding scenarios.
\end{itemize}

\begin{table*}
  \caption{\label{method-comparison}
    Comparison of different methods in terms of prefix sharing, dynamic dispatch,
    decision space, plan complexity, and Q-Bimodal support.
  }
  \centering
  \begin{tabular}{llllll}
    \hline
    \textbf{Method} & \textbf{Share Prefix} & \textbf{Dynamic Dispatch} & \textbf{Decision Space} & \textbf{Plan Complexity} & \textbf{Q-Bimodal} \\
    \hline
    Paged    & no  & no  & $O(1)$   & --              & -  \\
    FlashInfer     & yes & no  & $O(1)$   & --        & no  \\
    DeFT     & yes & no  & $O(1)$   & --              & no  \\
    FastTree & yes & yes & $O(2^K)$ & $O(K \log K)$   & no  \\
    Ours     & yes & yes & $O(1)$   & $O(K)$          & yes \\
    \hline
  \end{tabular}
\end{table*}


\section{Method}
\label{sec:method}

We propose a fixed-cardinality dispatcher for tree-structured attention.
The goal is to exploit KV sharing across branches while balancing between query tile padding overhead and extra decoding kernel launch overhead. At each step, the dispatcher summarizes the current batch, evaluates a small set of execution 
strategies, and launches the strategy with the lowest
predicted cost. 

\subsection{Problem Setup}

We consider a batch of prompts, each with $K$ active branches. At each
decoding step, every branch generates one query token. Branches may share a
common prefix, may share intermediate KV spans, and eventually diverge into
branch-specific tails.

\subsection{Dispatch Space}

Instead of searching over all nodes of the KV-sharing tree, we restrict the
runtime to four strategies:
\begin{equation}
\label{eq:strategy-set}
\mathcal{S}
=
\{s_{\rm pb}, s_{\rm 2f}, s_{\rm 2d}, s_{\rm 3d}\}.
\end{equation}
Here, $s_{\rm pb}$ is independent per-branch decoding, $s_{\rm 2f}$ is
two-level shared-prefix execution with a fused prefill-style tail, $s_{\rm 2d}$
is two-level shared-prefix execution with a decode-native tail, and
$s_{\rm 3d}$ is three-level shared-prefix execution in which the first two
shared levels are processed using fused prefill-style kernels, while the final
branch-specific level is processed using a decode-native kernel.

\subsection{Cost Model}

At every decoding step $t$ with work $W_t$, the dispatcher selects the strategy with minimum
predicted cost:
\begin{equation}
\label{eq:picker}
s_t^\star
=
\arg\min_{s\in\mathcal{S}}
C_s(W_t;\theta).
\end{equation}
where $\theta$ is a set of device-calibrated constants.
For fused shared-prefix strategies, the cost contains an effective bandwidth
term and a wave-quantized tile term:
\begin{equation}
\label{eq:fused-cost}
\begin{aligned}
C_{\rm fused}(s)
&=
B_s/\beta
+
\left\lceil
\frac{M_s^{\rm pref}}{N_{\rm SM}}
\right\rceil \tau
\\
&\quad
+
\left\lceil \log_2 D_s \right\rceil \tau_{\rm merge}.
\end{aligned}
\end{equation}

Here, $B_s = \sum_{\ell=1}^{D_s} b_\ell$ is the total KV bytes loaded across
the cascade levels of strategy $s$, $\beta$ is calibrated HBM bandwidth,
$M_s^{\rm pref}$ is the number of prefix tiles, $N_{\rm SM}$ is the number of
GPU SMs, and $\tau$ is calibrated time per tile. The final term models the parallel merge over the $D_s$ cascade levels, so a three-level strategy incurs two merge rounds while a two-level strategy incurs one.

For the decode-tail strategy, the prefix is costed as a fused shared-prefix
level, while the tail is costed with a decode-native model:
\begin{equation}
\label{eq:dectail-tail}
\begin{aligned}
C_{\rm dt}
&=
C_{\rm prefix}
+
\frac{B_{\rm tail}}{\beta}
+
\left\lceil
\frac{M_s^{\rm tail}}{N_{\rm SM}}
\right\rceil
\tau_{\rm tail}.
\end{aligned}
\end{equation}
Here, $B_{\rm tail}$ is tail KV traffic, $M_s^{\rm tail}$ is the number of
decode-tail work tiles, and $\tau_{\rm tail}$ is calibrated decode-kernel work. The calibration takes around 3 minutes on a H100 GPU per model. 

\subsection{Planning Complexity}

The dispatcher evaluates a constant number of strategies and caps the exposed
cascade depth at $D_{\max}=3$. Therefore, planning cost is
\begin{equation}
\label{eq:complexity}
O(|\mathcal{S}|\, K D_{\max})
=
O(K).
\end{equation}
This bound is independent of the number of nodes in the underlying KV-sharing
tree. Deeper sharing structure, if present, is folded into the branch-tail
term rather than introducing additional planner states. The result is a runtime that can dynamically choose both the exposed cascade
depth and the tail execution mode while retaining a fixed decision space.

\section{Experiments}
\subsection{Experiment Setup}

We evaluate our dispatcher on tree-structured decoding workloads using
Llama-3.2-1B, Llama-3.1-8B on 1 H100 GPU, and Llama-3-70B on 4 H100 GPU with tensor parallelism. We compare against Paged, FlashInfer's MLCA,
FastTree, and DeFT. We use the public released implementations of DeFT and FastTree, both of which report H100 results in their original papers, and use FlashInfer’s production-ready attention kernels for FlashInfer-based baselines. Our kernel is implemented using FlashInfer’s attention template, so it shares the same underlying attention-kernel infrastructure, additional details are in Appendix \ref{appendix:implementations}.

\subsection{Scenarios}
We evaluate four tree-structured decoding scenarios that cover different prefix lengths, batch sizes, and tree shapes. 
Here, $K$ denotes the number of leaves in the decoding tree.

\textbf{multi-doc-qa}~\cite{yang2018hotpotqa,10.1145/3701716.3715490} represents multi-document question answering workloads, where multiple questions are asked over several concatenated documents. This setting has a very long shared prefix, up to 80K tokens, a small batch size of $B=1$, and a large number of branches, $K=128$, and an output length of 256 tokens.

\textbf{multi-few-shot}~\cite{10.5555/3495724.3495883} represents few-shot prompting workloads, where the prompt contains multiple task demonstrations and the model generates outputs for multiple tasks. We use a relatively short shared prefix of 4K tokens, batch size $B=16$ for the 1B model and $B=4$ for the 8B and 70B models, $K=128$ branches, and an output length of 256 tokens.

\textbf{self-consistency}~\cite{wang2023selfconsistencyimproveschainthought} represents parallel
  sampling workloads, where a single prompt is expanded into many independent
  candidate continuations that are aggregated after decoding---for example,
  majority voting over chain-of-thought samples, or best-of-$N$ selection with a
  reward model or verifier. Unlike the multi-query scenarios above, the $K$
  branches all share one common prefix and never re-branch, forming a
  \emph{static} two-level tree: a single shared context fanning out into $K$
  parallel tails. This setting has a moderate prefix of 16K tokens, a small batch
  size of $B=1$, and $K=16$ branches. To do a  long decode stress case, we output a length of 12k tokens.
  
\textbf{beam-search}~\cite{vijayakumar2018diversebeamsearchdecoding} represents a dynamic decoding workload in which the output expands into a tree according to accumulated logit scores. The goal is to explore multiple hypotheses and find a higher-scoring continuation. Unlike the previous two scenarios, the tree structure in beam search is multi-level and difficult to predict in advance, making it the most challenging setting for tree attention. We use a 40K-token prefix, a moderate batch size of $B=16$, $K=64$ branches, and an output length of 256 tokens.

\subsection{Results}
Figure~\ref{fig:e2e_bar} shows the end-to-end model execution performance across all evaluated scenarios. Our method achieves the best performance in every setting and for all evaluated model sizes. Compared with PagedAttention, our method delivers substantial speedups, ranging from $1.02\times$ to $6.40\times$, demonstrating the benefit of exploiting shared prefixes and tree structure during decoding.

FastTree and MLCA are the strongest competing baselines, with the better method varying across settings. Even when comparing against the stronger of these two baselines for each configuration, Our method consistently delivers speedups across all settings, ranging from 1.11× to 2.73×. This demonstrates the robustness and effectiveness of the Q-Bimodal design.

The speedup varies with workload structure and model configuration. The 1B and 8B models, both with a GQA ratio of 4, achieve higher speedups than the 70B model, which has a GQA ratio of 8. This is expected, since models with a higher GQA ratio have a smaller per-query KV footprint. Workloads with longer shared prefixes, larger batch sizes, or more branches expose greater redundancy in PagedAttention’s memory traffic, enabling our method to deliver larger performance gains. 

Table~\ref{tab:time-breakdown} further breaks down execution time across planning and forward computation. The dispatcher keeps planning overhead low, while the Q-Bimodal kernel substantially reduces forward time compared with all baselines. These results show that the overall speedup comes primarily from the kernel-level improvements.

\begin{figure}[t]
\centering
    \includegraphics[width=0.8\linewidth]{llama_speedup_bars.pdf}
    \caption{
    TPOT speedup over PagedAttention.
    PagedAttention is normalized to $1.0\times$ for each setting, and higher is better.
    F1 is multi-document QA, F2 is multi-few-shot, F3 is beam-search, and F4 is self-consistency.}
    \label{fig:e2e_bar}
\end{figure}



\begin{table}[t]
\centering
\caption{\label{tab:time-breakdown}
Timing breakdown for Llama-3.1-8B on \texttt{beam-search}. 
Build metadata and decision time capture CPU overhead, with decision time corresponding to kernel-strategy selection. 
Forward time is the model forward-pass GPU computation time. 
All times are reported in milliseconds.
}
\resizebox{\columnwidth}{!}{%
\begin{tabular}{lrrr}
  \hline
  \textbf{Method} &
  \textbf{Build Metadata} &
  \textbf{Decision} &
  \textbf{Forward} \\
  \hline
  DeFT     & 1{,}533 & 0   & 21{,}382 \\
  MLCA     & 1{,}417 & 0   & 12{,}493 \\
  Ours     & 1{,}167 & 93  & 7{,}717  \\
  FastTree & 1{,}962 & 786 & 19{,}083 \\
  Paged    & 960     & 0   & 37{,}215 \\
  \hline
\end{tabular}%
}
\end{table}

\subsection{Dispatcher Analysis}
We analyze the execution plans selected by the runtime dispatcher. For each workload, we compare the adaptive dispatch policy against runs in which a single dispatch strategy is forced for the entire decode. We use the fastest forced strategy as a static oracle. Although the true optimum may vary across decode steps, exhaustively searching over all per-step strategy sequences is prohibitively expensive. The best static policy therefore provides a practical and strong oracle baseline.

Table~\ref{tab:picker_diagnosis} reports the dispatcher diagnosis results.
We omit $s_{\rm 3d}$ from the table because it is never faster than $s_{\rm 2d}$ in our experiments. This is notable because beam-search can expose more than three sharing levels. However, the additional intermediate levels tend to be short and have too few descendants to offset the extra merge and launch overheads, making deeper sharing unprofitable.

The self-consistency benchmark highlights the value of dynamic dispatch. Early in decoding, when there is limited tail work, the dispatcher selects $s_{\rm 2f}$ to avoid the overhead of a separate decode-tail kernel. As the tails grow longer, it switches to $s_{\rm 2d}$, where decode-native tail execution becomes more efficient. This adaptive behavior allows the dispatcher to outperform both static strategies for the 1B and 8B models, demonstrating that no single static execution strategy is uniformly optimal across all decode steps.

Overall, the dispatcher closely tracks the best available strategy with little runtime overhead. In most settings, it matches the fastest static policy within a small margin, while in workloads whose optimal strategy changes over time, it can exceed the best static baseline by switching strategies during decoding.


\begin{table}[t]
\centering
\small
\resizebox{\columnwidth}{!}{
\begin{tabular}{llrrrl}
\toprule
Scenario & Model & $s_{\rm 2f}$-only & $s_{\rm 2d}$-only & Dispatcher &  $s_{\rm 2f}/s_{\rm 2d}$ \\
\midrule
\multirow{3}{*}{F1}
& 1B  & 9{,}929  & 2{,}241  & 2{,}243  & $s_{\rm 2d}$ \\
& 8B  & 29{,}091 & 6{,}609  & 6{,}619  & $s_{\rm 2d}$ \\
& 70B & 52{,}139 & 11{,}317 & 11{,}348 & $s_{\rm 2d}$ \\
\midrule
\multirow{3}{*}{F2}
& 1B  & 9{,}152  & 7{,}757  & 7{,}791  & $s_{\rm 2d}$ \\
& 8B  & 8{,}471  & 6{,}829  & 6{,}850  & $s_{\rm 2d}$ \\
& 70B & 16{,}398 & 14{,}330 & 14{,}340 & $s_{\rm 2d}$ \\
\midrule
\multirow{3}{*}{F3}
& 1B  & 9{,}862  & 8{,}731  & 8{,}860  & $s_{\rm 2d}$ \\
& 8B  & 12{,}971 & 9{,}871  & 9{,}909  & $s_{\rm 2d}$ \\
& 70B & 26{,}309 & 18{,}054 & 18{,}158 & $s_{\rm 2d}$ \\
\midrule
\multirow{3}{*}{F4}
& 1B  & 62{,}819  & 61{,}451  & 59{,}637  & 35/65 \\
& 8B  & 169{,}475 & 160{,}849 & 130{,}132 & 17/83 \\
& 70B & 377{,}826 & 406{,}988 & 379{,}352 & 72/28 \\
\bottomrule
\end{tabular}
}
\caption{
Dispatcher diagnosis across models and scenarios. Lower is better.
The $s_{\rm 2f}$-only and $s_{\rm 2d}$-only columns report forced static execution with a single strategy. F1 = multi-doc-qa,
F2 = multi-few-shot, F3 = beam-search, and F4 = self-consistency.
}
\label{tab:picker_diagnosis}
\end{table}


\section{Conclusion}

We presented Q-Bimodal Tree Attention, a fixed-cardinality dispatcher for
efficient tree-structured LLM decoding. By separating shared-prefix execution
from branch-specific tail execution, our method reduces redundant KV reads and
tail-side query padding while keeping planning cost linear in the number of
branches. Experiments across diverse decoding workloads show consistent
speedups over PagedAttention, fixed shared-prefix methods, DeFT, MLCA, and
FastTree. These results show that Q-Bimodal execution is the key to efficient
tree-structured decoding, as it adapts between fused shared-prefix attention
and decode-native tail attention according to the workload shape.

\section*{Limitations}

Our cost model is a simple estimator rather than an exact runtime predictor. It is
designed to provide a good-enough ranking of execution strategies with low
decision overhead, not to precisely estimate the latency of each strategy.

Our dispatcher also covers only a fixed set of strategies. One can certainly
construct tree shapes for which these fixed strategies are not optimal.
However, for recent tree-structured decoding applications, this restricted
strategy set is already sufficient and performs well in practice.


\bibliography{custom}


\appendix
\label{sec:appendix}


\section{Baseline and Implementations}
  \label{appendix:implementations}
  \paragraph{Hardware and software.} All methods run inside the same
  driver with the same paged KV allocator (page size 16, refcounted
  copy-on-write) and the same standalone Llama implementation, so
  weight loading, RMSNorm, RoPE, and projections are identical across
  rows methods. Code is Python~3.11 on
  PyTorch~2.9 with the FlashInfer attention library. Methods differ only
  in the per-step attention call.

  \paragraph{FastTree~\cite{pan2025fasttree}.} Public artifact,
  integrated into our driver so allocator, model, and batching match
  the other rows. Plan-time builds a page-level radix tree and runs
  the FastTree heuristic to score per-vnode
  \texttt{(SplitQ, SplitK)} against offline-autotuned thresholds. To
  make the comparison fair we applied an LCA-prefix-skip in the radix
  walker, a numpy-vectorized page-to-slot expansion, and a batched
  host-to-device transfer of the per-step metadata; the heuristic loop
  remains the dominant plan-time cost.

  \paragraph{DeFT~\cite{yao2025deft}.} Vendored kernel (DeFT pins
  \texttt{torch}=2.5.1 and cannot be co-installed). We reuse the
  page-level radix tree from the FastTree integration and convert it
  to DeFT's flat-array metadata; nodes whose reader count exceeds
  $\mathtt{BLOCK\_M}=32$ are split into
  $\lceil n_{\text{readers}}/32 \rceil$ virtual entries sharing the
  same KV range. A per-prompt plan cache amortizes metadata builds
  across non-forking steps.



\section{Dispatch Space Ablation}

We next ask whether additional dispatch strategies for deeper sharing
levels are necessary. This ablation uses the Llama-3.1-8B long-decode
beam-search setting with $K=32$, $L_p=8192$, $B=8$, and
does 2000 decode steps. We force the dispatcher to use deeper variants and measure only decode steps 1500--2000, where the
beam-search tree is most likely to contain longer intermediate shared
prefixes that deeper strategies could exploit.

\begin{table}[H]
\centering
\caption{
Depth ablation on Llama-3.1-8B with TP=2. Forward time is measured
over decode steps 1500--2000. Improvement is relative to the
corresponding two-level variant within each kernel family.
}
\label{tab:llama8b_depth_ablation}
\begin{tabular}{@{}lrr@{}}
\toprule
Strategy & Fwd. time (ms) & Imp. vs 2L \\
\midrule
  2L\_DECTAIL & 10{,}739 & -- \\                             
  3L\_DECTAIL & 11{,}651 & +8.5\% \\                          
  4L\_DECTAIL & 10{,}731 & -0.1\% \\                          
  5L\_DECTAIL & 10{,}999 & +2.4\% \\                                                 
\bottomrule                       
\end{tabular} 
\end{table}

The results in Table~\ref{tab:llama8b_depth_ablation} show that increasing the number of fused cascade
levels does not materially improve forward time, as more levels can only share shorter prefixes and brings more merge overhead.

\section{Query-Tile Padding Overhead of $s_{2f}$}
\label{appendix:padding}

The only structural difference between $s_{2f}$ and $s_{2d}$ is how the
branch-specific tail level is executed. $s_{2f}$ folds the tail into the
same fused prefill-style kernel that processes the shared prefix, whereas
$s_{2d}$ runs a decode-native kernel for the tail and merges it with the
prefix through an online-softmax bridge. This appendix isolates the cost
of that single choice, which is the query-tile padding paid by $s_{2f}$.

\paragraph{Source of the overhead.}
A prefill-style kernel processes queries in tiles of a fixed width $T$
(the kernel's \texttt{CTA\_TILE\_Q}). As in the per-tile cost model of
Section~\ref{sec:method}, a tile issues the full $\mathbf{Q}\mathbf{K}^\top$
and $\mathbf{P}\mathbf{V}$ MMA chains over its entire KV span regardless of
how many query rows it actually carries; unused rows are merely masked at
writeback, so the cost is set by $T$ rather than by the real row count. At
a decode step, each branch contributes only $G$ real query rows --- one
generated token folded over its GQA group of size
$G = n_{\rm qo}/n_{\rm kv}$ --- yet $s_{2f}$ pads every tail tile to width
$T$. The MMA utilization of a tail tile is therefore
\begin{equation}
\label{eq:pad-util}
U = \frac{G}{T},
\qquad
\text{wasted fraction} = 1 - \frac{G}{T},
\end{equation}
and the wasted work scales with the tail KV length, so the penalty grows
as the tails lengthen during decoding. The decode-native tail of $s_{2d}$
processes exactly the $G$ rows per branch (\texttt{CTA\_Q}$=1$ with GQA
fold-in) and incurs no query padding.

\paragraph{Padding fraction.}
Because $U=G/T$ is fixed by the model's GQA ratio and the launched tile
width, the per-tile padding fraction is exact. $s_{2f}$ uses a single
fused pool, so its tail rows are launched at the same large tile width
$T\in\{64,128\}$ as the shared prefix; the smallest prefill tile a fused
shared-prefix engine would use for tails is $T=16$. Table~\ref{tab:padding_fraction}
lists the resulting waste. Even at the most favorable $T=16$, $s_{2f}$
discards $50$--$75\%$ of its tail MMA cycles; at the $T\in\{64,128\}$
widths it actually launches, the waste exceeds $87\%$. We confirmed this
empirically by instrumenting the launched tiles over a 127-step decode on
Llama-3.2-1B ($K=32$, $B=32$, $L_p=8192$): a prefill-style tail tile at
$T=16$ records a $75\%$ query-padding fraction every step --- matching
$1-G/T$ exactly --- while the decode-native tail records $0\%$ by
construction.

\begin{table}[H]
\centering
\caption{
Tail-tile query-padding fraction $1-G/T$ for the fused prefill-style tail
of $s_{2f}$, by GQA group size $G$ and launched tile width $T$. The
decode-native tail of $s_{2d}$ has $0\%$ padding for every entry.
}
\label{tab:padding_fraction}
\begin{tabular}{@{}llrrr@{}}
\toprule
Model & $G$ & $T=16$ & $T=64$ & $T=128$ \\
\midrule
1B  & 4 & 75.0\% & 93.75\% & 96.88\% \\
8B  & 4 & 75.0\% & 93.75\% & 96.88\% \\
70B & 8 & 50.0\% & 87.50\% & 93.75\% \\
\bottomrule
\end{tabular}
\end{table}

\paragraph{End-to-end effect.}
The padding waste is most exposed when the branch tails are long and the
shared prefix is comparatively cheap, so that the tail MMA dominates the
step. The multi-doc-qa scenario (F1 in Table~\ref{tab:picker_diagnosis})
is exactly this regime: the $s_{2f}$-only column is $4.4$--$4.6\times$
slower than $s_{2d}$-only across all three models. In the opposite regime
--- short tails with a large shared prefix --- the padding is amortized
against the prefix read that both strategies share, and the gap collapses
(e.g.\ F2 on the 8B model, where $s_{2f}$ is within $1.24\times$ of
$s_{2d}$). This crossover is precisely what the dispatcher exploits: it
launches $s_{2f}$ while the tails are short and the extra decode-kernel
launch of $s_{2d}$ is not yet worth paying, then switches to $s_{2d}$ once
the tails grow long enough that the padding tax dominates.

\section{Calibration}
  \label{appendix:calibration}

  The cost model in Section~\ref{sec:method} is parameterized by a small set
  of device-calibrated constants
  $\theta = \{\beta, \tau, \tau_{\rm merge}, \tau_{\rm tail}$.
  We do \emph{not} re-tune the attention kernels themselves: FlashInfer's
  JIT already specializes register allocation, shared-memory layout, and
  instruction selection per SM, so calibration only measures the
  behaviors the JIT cannot expose at decision time. Each constant is
  recovered by a short microbenchmark run once at engine initialization.
  The number of SMs $N_{\rm SM}$ is read directly from the device
  properties and is not timed. The full sweep takes roughly three minutes
  on an H100.

  \paragraph{Bandwidth $\beta$.}
  We measure effective HBM bandwidth with a device-to-device copy of a
  256\,MiB slab, averaged over several runs, and report $\beta$ in bytes
  per microsecond. This anchors every bandwidth term $B_s/\beta$ in
  Equations~\ref{eq:fused-cost} and~\ref{eq:dectail-tail}; on H100 it
  recovers roughly 1.4\,TB/s.

  \paragraph{Per-tile prefill cost $\tau$.}
  We time a single fused prefill call with $T$ query rows against a
  fixed-length KV span, swept over the tile widths the kernel
  JIT-compiles, $T \in \{16, 64, 128\}$. This yields the per-tile term
  used by the wave-quantized prefix cost.
  
\paragraph{Decode-tail cost $\tau_{\rm tail}$.}
  We time a paged-decode call at one query per branch
  ($\mathrm{CTA}_Q = 1$) against KV spans of varying length, and fit the
  per-token decode work from how latency grows with the KV tokens loaded.
  This yields the $\tau_{\rm tail}$ term used by the wave-quantized
  decode-tail cost in Equation~\ref{eq:dectail-tail}.
  
\paragraph{Merge cost $\tau_{\rm merge}$.}
  The per-round merge cost $\tau_{\rm merge}$ in Equation~\ref{eq:fused-cost}
  is calibrated as a single constant, measured from the latency of one
  \texttt{merge\_state\_in\_place} call. A depth-$D_s$ cascade combines its
  levels in $\lceil \log_2 D_s \rceil$ parallel merge rounds along the
  critical path, matching the final term of
  Equation~\ref{eq:fused-cost}.

  \section{Reproduction}
  \label{appendix:reproduction}

  \paragraph{Models and precision.}
  Llama-3.2-1B-Instruct and Llama-3.1-8B-Instruct are served with
  \texttt{bf16} weights and a \texttt{bf16} KV cache on a single H100
  (no tensor parallelism). Llama-3-70B-Instruct is served with
  \texttt{bf16} weights and a \texttt{bf16} KV cache under tensor
  parallelism (TP$=4$) across four H100s. All checkpoints are the public
  instruction-tuned releases.

  \paragraph{Workload data.}
  The four scenarios are built from two public datasets, GSM8K and
  HotpotQA~\cite{yang2018hotpotqa}, using real text rather than synthetic
  tokens. \emph{multi-doc-qa} (F1) concatenates HotpotQA documents into
  an $\sim$80K-token shared prefix and appends $K{=}128$ distinct
  HotpotQA questions. \emph{multi-few-shot} (F2) and
  \emph{self-consistency} (F4) use GSM8K few-shot exemplars as the shared
  prefix; F4 pads the exemplars-plus-question context to the target
  prefix length and fans it into $K$ independent samples.
  \emph{beam-search} (F3) uses HotpotQA prompts as the prefix and grows
  the $B{\times}K$ beam tree online from accumulated logit scores, so its
  multi-level sharing structure is data-dependent and not fixed in
  advance. The GSM8K test split is downloaded on first use; the HotpotQA
  multi-question file is built by a released preprocessing script.

  \paragraph{Timing protocol.}
  For every (scenario, method) cell we run a warmup iteration, discard
  it, and report the median over the measured repeats; the CUDA cache is
  cleared and garbage collection is run between repeats. We report total
  decode time over all steps, split into \emph{plan} time (CPU-side
  metadata build plus strategy decision) and \emph{forward} time (GPU
  compute). F1--F3 decode 256 tokens per branch; F4 is the long-decode
  stress cell at 12{,}000 tokens. The speedups in
  Figure~\ref{fig:e2e_bar} are time-per-output-token (TPOT) relative to
  PagedAttention in the same cell.

  \paragraph{Calibration and determinism.}
  The cost-model constants $\theta$ are recalibrated once per model at
  engine initialization (Appendix~\ref{appendix:calibration}); the sweep
  takes roughly three minutes. Every method shares the same paged
  allocator, model implementation, and sampler, and all methods produce
  token-identical greedy outputs, so the reported differences reflect
  only attention-execution time.
  
\section{Artifact Licenses and Terms of Use}
\label{appendix:licenses}

We use publicly available software, models, datasets, and baseline artifacts
whose licenses or terms permit research use. Our released artifact will contain
our implementation, benchmark scripts, configuration files, and instructions
needed to reproduce the experiments. Unless otherwise stated in the artifact
repository, we release our own code under the Apache-2.0 license. The artifact
does not redistribute third-party model weights, datasets, CUDA/NVIDIA software,
or vendored dependencies whose licenses require users to obtain them separately.

\paragraph{Software dependencies.}
Our implementation is written in Python and uses PyTorch and FlashInfer as the
main software dependencies. Python is distributed under the Python Software
Foundation License Version 2. PyTorch is distributed under a permissive
BSD-style license. FlashInfer is distributed under the Apache-2.0 license. For
PagedAttention baselines, we rely on the public vLLM implementation, which is
also distributed under the Apache-2.0 license. CUDA and NVIDIA GPU software are
used subject to the NVIDIA CUDA Toolkit End User License Agreement and
applicable NVIDIA driver/software license terms.

\paragraph{Models.}
The experiments use Llama-3.2-1B, Llama-3.1-8B, and Llama-3-70B. These models
are governed respectively by the Llama 3.2 Community License Agreement, the
Llama 3.1 Community License Agreement, and the Meta Llama 3 Community License
Agreement, together with their Acceptable Use Policies. We do not redistribute
the model weights in our artifact. Users who reproduce the experiments must
obtain the models from the official Meta or Hugging Face distribution channels
and agree to the applicable Llama license terms.

\paragraph{Datasets and workloads.}
The multi-document QA workload uses HotpotQA. The HotpotQA dataset is
distributed under the CC BY-SA 4.0 license, while the accompanying HotpotQA code
is distributed under the Apache-2.0 license. Other workloads in our evaluation,
including beam search, self-consistency, and multi-few-shot prompting, are
synthetic or procedure-defined decoding workloads derived from standard
generation settings rather than newly collected human-subject data.

\paragraph{Third-party research artifacts.}
We compare against public implementations of DeFT and FastTree. DeFT is released
under the MIT license. The FastTree artifact for tree-structured LLM inference
is released under the Apache-2.0 license. We use these artifacts only for
benchmarking and baseline comparison, and we preserve their copyright notices,
license files, and citation requirements. If our artifact includes wrapper code
or integration scripts for these baselines, we document the original upstream
repositories and licenses.

\paragraph{Use and distribution.}
The intended use of our artifact is academic research and reproducibility of the
results reported in this paper. Users are responsible for complying with the
licenses and terms of all third-party components, including software libraries,
model weights, datasets, CUDA/NVIDIA software, and baseline artifacts. No
additional restrictions are imposed by us beyond the license selected for our
own released code and the terms of the third-party artifacts listed above.

\end{document}

